% =====================================================================
%  Robert Hines, "Topological K-theory" (2016-12-15) —— 中文译注本
%  本文件为导言区（Lead 独占；译者只读，不得修改）
%
%  编译：LuaLaTeX（本机 xdvipdfmx 不稳，统一用 lualatex，编译 3 遍）
%        powershell -ExecutionPolicy Bypass -File build.ps1
%
%  约定：
%   * 原文无编号公式（实测 0 条），故本译本不使用 \tag；原文的编号定理/定义
%     一律手工给号：\begin{thm}{1} … \end{thm}（同理 defn/prop/lem）。
%   * 原文的展示式全部逐符号重排，不用截图。
%   * % [Tutor Note: ...] 为译者批注，不进入排版，汇总见《译注汇总.md》。
% =====================================================================
\documentclass[11pt,a4paper]{article}

\usepackage[UTF8]{ctex}
\usepackage{amsmath,amssymb,amsthm,mathtools}
\usepackage{bm}
\usepackage{graphicx}
\usepackage[a4paper,left=2.6cm,right=2.6cm,top=2.6cm,bottom=2.7cm]{geometry}
\usepackage[dvipsnames]{xcolor}
\usepackage{enumitem}
\usepackage{booktabs}
\usepackage{array}
\usepackage{multicol}
\usepackage{fancyhdr}
\usepackage{caption}
\usepackage{tikz}
\usetikzlibrary{cd,arrows.meta,positioning}
\usepackage[hidelinks,unicode]{hyperref}
\hypersetup{pdftitle={拓扑 K 理论（中译注本）},
  pdfauthor={Robert Hines; 中译/批注: 本译本},
  pdfsubject={Hines, Topological K-theory, 2016-12-15}}

\captionsetup{font=small,labelsep=quad}
\setlength{\emergencystretch}{3em}
\allowdisplaybreaks[1]

% ---------- 页眉 ----------
\pagestyle{fancy}
\fancyhf{}
\fancyhead[L]{\small\itshape Topological K-theory \textnormal{(Hines, 2016)} 中译注本}
\fancyhead[R]{\small\thepage}
\renewcommand{\headrulewidth}{0.4pt}

% ---------- 定理类环境：编号一律手工给出（与原文逐条对应） ----------
% 中文版体例：粗体 “定义 1.” + 直立体正文（中文不用斜体；原文的斜体是西文排版习惯）。
\newenvironment{defn}[1]{\par\medskip\noindent\textbf{定义 #1.}\ }{\par\medskip}
\newenvironment{prop}[1]{\par\medskip\noindent\textbf{命题 #1.}\ }{\par\medskip}
\newenvironment{lem}[1]{\par\medskip\noindent\textbf{引理 #1.}\ }{\par\medskip}
\newenvironment{thm}[1]{\par\medskip\noindent\textbf{定理 #1.}\ }{\par\medskip}
\renewenvironment{proof}[1][证明]{\par\medskip\noindent\textbf{#1.}\ }{\hfill$\square$\par\medskip}

% 原文的行内粗体小标题：独立成段、粗体、不缩进
\newcommand{\para}[1]{\par\medskip\noindent\textbf{#1}\quad\ignorespaces}

% 原文图/表标题（本稿实测无图无表；留宏以备后用）
\newcommand{\figtitle}[2]{\caption*{\small\textbf{图 #1：}#2}}
\newcommand{\tabtitle}[2]{\caption*{\small\textbf{表 #1：}#2}}

% 参考文献条目：\refitem{键}{内容}
\newcommand{\refitem}[2]{\par\hangindent=1.5em\hangafter=1\noindent [#1]~#2}

% ---------- 常用记号（只能用这里定义过的宏，不要自造 \newcommand） ----------
\newcommand{\R}{\mathbb{R}}
\newcommand{\C}{\mathbb{C}}
\newcommand{\Z}{\mathbb{Z}}
\newcommand{\N}{\mathbb{N}}
\newcommand{\Q}{\mathbb{Q}}
\newcommand{\GL}{\operatorname{GL}}
\newcommand{\Hom}{\operatorname{Hom}}
\newcommand{\End}{\operatorname{End}}
\newcommand{\Aut}{\operatorname{Aut}}
\newcommand{\Vect}{\mathrm{Vect}}
\newcommand{\dd}{\mathrm{d}}
\newcommand{\id}{\mathrm{id}}
\newcommand{\rk}{\operatorname{rk}}
\newcommand{\ie}{\textit{i.e.}}
\newcommand{\eg}{\textit{e.g.}}
\newcommand{\cf}{\textit{cf.}}
